\chapter{Código-Fonte Selecionado}
\lhead{Apêndice \thechapter: Código-Fonte Selecionado}\rhead{\thepage}

Este apêndice apresenta o código-fonte das interfaces (arquivos de cabeçalho \texttt{.h}) das principais classes desenvolvidas no projeto. A escolha por apresentar as interfaces visa a oferecer uma visão clara e concisa da arquitetura orientada a objeto, das responsabilidades de cada classe e de suas inter-relações, sem o detalhamento excessivo dos arquivos de implementação.

\section{Classe CSimulador}
Responsável pela orquestração geral do sistema, gerenciando o fluxo de execução e a interação entre os componentes.
\begin{verbatim}
#ifndef CPECAPPLICATION_H
#define CPECAPPLICATION_H

#include "CInput.h"
#include "CElasticCalculator.h"
#include "CMineralDatabase.h"
#include "CReservoirData.h"
#include "CPlotManager.h"

class CSimulador {
public:
    CSimulador();
    void Run();

private:
    CInput m_inputManager;
    CMineralDatabase m_database;
    CElasticCalculator m_calculator;
    CReservoirData m_reservoirData;
    CPlotManager m_plotManager;
    bool m_running;

    void Menu();
    void Initialize();
    void Finalize();
};

#endif
\end{verbatim}

\section{Classe CElasticCalculator}
Encapsula a lógica de cálculo dos módulos elásticos, utilizando composição para delegar os cálculos específicos.
\begin{verbatim}
#ifndef CELASTICCALCULATOR_H
#define CELASTICCALCULATOR_H

#include "CMineralDatabase.h"
#include "CAmostra.h"
#include "CBulkCalculator.h"
#include "CShearCalculator.h"

class CElasticCalculator {
public:
    struct Result {
        float bulkVoigt, bulkReuss, bulkHill;
        float shearVoigt, shearReuss, shearHill;
    };

    explicit CElasticCalculator(const CMineralDatabase& db);
    Result Calculate(const CAmostra& amostra) const;

private:
    CBulkCalculator m_bulkCalc;
    CShearCalculator m_shearCalc;
};

#endif // CELASTICCALCULATOR_H
\end{verbatim}

\section{Classe CAmostra}
Representa a entidade de dados de uma amostra de rocha, encapsulando todas as suas propriedades.
\begin{verbatim}
#ifndef CAMOSTRA_H
#define CAMOSTRA_H

#include <string>
#include <map>

class CAmostra {
public:
    CAmostra();
    CAmostra(const std::string& reservatorio, const std::string& nome,
             float profundidade, float porosidade,
             const std::map<std::string, float>& minerais);

    const std::string& NomeReservatorio() const;
    const std::string& NomeAmostra() const;
    float Profundidade() const;
    float Porosidade() const;
    const std::map<std::string, float>& Minerais() const;

private:
    std::string m_reservatorio;
    std::string m_nome;
    float m_profundidade;
    float m_porosidade;
    std::map<std::string, float> m_minerais;
};

#endif
\end{verbatim}